\setlength{\parindent}{0pt}

\section{Matemática}

\subsection{Probabilidade}

\textbf{Bayes:} Atualiza a probabilidade de $A$ dado que $B$ ocorreu.
\[P(A\mid B)=\frac{P(B\mid A)\,P(A)}{P(B)}\]
Quando $P(B)$ não é conhecida diretamente, expanda pelo denominador via lei da probabilidade total ($A_j$ são hipóteses mutuamente exclusivas e exaustivas):
\[P(A_i\mid B)=\frac{P(B\mid A_i)\,P(A_i)}{\sum_j P(B\mid A_j)\,P(A_j)}\]

\textbf{Esperança:} $E[X]=\sum x\,P(X{=}x)$

Linearidade: $E[aX+bY]=aE[X]+bE[Y]$

Se independentes: $E[XY]=E[X]E[Y]$

\textbf{Variância:} $\text{Var}(X)=E[X^2]-E[X]^2$; 

$\text{Var}(aX+b)=a^2\text{Var}(X)$

Se independentes: $\text{Var}(X+Y)=\text{Var}(X)+\text{Var}(Y)$

\vspace{4pt}

\textbf{Markov:} Limite superior fraco para $X\geq0$; usa só a esperança.
\[P[X\geq\lambda E[X]]\leq1/\lambda\]

\textbf{Chebyshev:} Limita a probabilidade de $X$ se desviar muito da média; usa esperança e variância.
\[P[|X-E[X]|\geq\lambda\sigma]\leq1/\lambda^2\]

\noindent\rule{\linewidth}{0.2pt}

\textbf{Binomial} $B(n,p)$: $n$ ensaios de Bernoulli independentes com prob.\ $p$.
\[P(X=k)=\binom{n}{k}p^k(1-p)^{n-k}\]
\[E[X]=np, \quad \text{Var}=np(1-p)\]

\noindent\rule{\linewidth}{0.2pt}

\textbf{Geométrica} $\text{Geom}(p)$: número de tentativas até o primeiro sucesso ($q=1-p$).
\[P(X=k)=pq^{k-1}, \quad E[X]=\tfrac{1}{p}, \quad \text{Var}=\tfrac{q}{p^2}\]

\noindent\rule{\linewidth}{0.2pt}

\textbf{Hipergeométrica} $H(N,K,n)$: $N$ itens, $K$ marcados, $n$ sorteados sem reposição.
\[P(X=k)=\frac{\binom{K}{k}\binom{N-K}{n-k}}{\binom{N}{n}}, \quad E[X]=\frac{nK}{N}\]

\noindent\rule{\linewidth}{0.2pt}

\textbf{Poisson} $\text{Pois}(\lambda)$: aproximação de $B(n,p)$ quando $n\to\infty$, $np=\lambda$.
\[P(X=k)=\frac{e^{-\lambda}\lambda^k}{k!}, \quad E[X]=\text{Var}(X)=\lambda\]

\noindent\rule{\linewidth}{0.2pt}

\textbf{Normal} $\mathcal{N}(\mu,\sigma^2)$: $p(x)=\frac{1}{\sqrt{2\pi}\,\sigma}e^{-(x-\mu)^2/2\sigma^2}$, $E[X]=\mu$, $\text{Var}=\sigma^2$.

\noindent\rule{\linewidth}{0.2pt}

\textbf{Coupon collector:} $n$ tipos equiprováveis; número esperado de sorteios $=nH_n$.

\textbf{Paradoxo do aniversário:} $\approx1.18\sqrt{M}$ amostras de $[M]$ para colisão com prob $\geq50\%$.

\subsection{Trigonometria}

$\sin^2x+\cos^2x=1$, $1+\tan^2x=\sec^2x$, $1+\cot^2x=\csc^2x$

$\sin(x\pm y)=\sin x\cos y\pm\cos x\sin y$

$\cos(x\pm y)=\cos x\cos y\mp\sin x\sin y$

$\tan(x\pm y)=\dfrac{\tan x\pm\tan y}{1\mp\tan x\tan y}$

$\sin2x=2\sin x\cos x$, $\cos2x=\cos^2x-\sin^2x=2\cos^2x-1=1-2\sin^2x$

$e^{ix}=\cos x+i\sin x$, $e^{i\pi}=-1$

\textbf{Lei dos senos:} $\dfrac{a}{\sin A}=\dfrac{b}{\sin B}=\dfrac{c}{\sin C}=2R$

\textbf{Lei dos cossenos:} $c^2=a^2+b^2-2ab\cos C$

\textbf{Área do triângulo:}

\begin{center}
\begin{tikzpicture}[scale=1.1, font=\small]
  \coordinate (A) at (0,0);
  \coordinate (B) at (3,0);
  \coordinate (C) at (1,1.8);
  \draw (A) -- node[below] {$c$} (B)
            -- node[right] {$a$} (C)
            -- node[left]  {$b$} cycle;
  \node[below left]  at (A) {$A$};
  \node[below right] at (B) {$B$};
  \node[above]       at (C) {$C$};
\end{tikzpicture}
\end{center}

\[\tfrac{1}{2}ab\sin C\]
\[\sqrt{s(s-a)(s-b)(s-c)},\ s=\tfrac{a+b+c}{2}\]
\[\tfrac{c^2\sin A\sin B}{2\sin C},\ C=\pi-A-B\]

\textbf{Circumraio:} $R=abc/(4A)$. \quad \textbf{Inraio:} $r=A/s$.

\subsection{Geometria}

\textbf{Shoelace:} $A=\frac12\left|\sum_i(x_iy_{i+1}-x_{i+1}y_i)\right|$

\textbf{Pick:} $A=I+B/2-1$ ($I$ = pontos interiores, $B$ = pontos na borda)

\textbf{Área por coordenadas:} $\dfrac{1}{2}\left|\det\begin{pmatrix}x_1-x_3&x_2-x_3\\y_1-y_3&y_2-y_3\end{pmatrix}\right|$

\textbf{Esferas/círculos:} $A=\pi r^2$, $V=\tfrac43\pi r^3$.

\textbf{Produto vetorial:} $\vec{u}\times\vec{v}=(u_yv_z-u_zv_y,\;u_zv_x-u_xv_z,\;u_xv_y-u_yv_x)$

Em 2D: $u_xv_y-u_yv_x$ ($>0$ anti-horário, $=0$ colinear; módulo = área do paralelogramo)

\textbf{Dist. ponto--ponto:} $d=\sqrt{(x_1-x_2)^2+(y_1-y_2)^2}$

(3D: inclui $(z_1-z_2)^2$)

\textbf{Dist. ponto--reta} ($ax+by+c=0$): $d=\dfrac{|ax_0+by_0+c|}{\sqrt{a^2+b^2}}$

\textbf{Reta} por $(x_0,y_0)$ e $(x_1,y_1)$: $\left|\begin{smallmatrix}x&y&1\\x_0&y_0&1\\x_1&y_1&1\end{smallmatrix}\right|=0$

\textbf{Rotação 2D:} $\begin{pmatrix}x'\\y'\end{pmatrix}=\begin{pmatrix}\cos\theta&-\sin\theta\\\sin\theta&\cos\theta\end{pmatrix}\begin{pmatrix}x\\y\end{pmatrix}$

\textbf{Rotação 3D} em torno de $x$, $y$, $z$ por $\theta$:
\[R_x=\begin{pmatrix}1&0&0\\0&c&-s\\0&s&c\end{pmatrix},\quad
R_y=\begin{pmatrix}c&0&s\\0&1&0\\-s&0&c\end{pmatrix}\]
\[R_z=\begin{pmatrix}c&-s&0\\s&c&0\\0&0&1\end{pmatrix}\]
onde $c=\cos\theta$, $s=\sin\theta$.

